\documentclass[11pt,a4paper]{article}

\usepackage[T1]{fontenc}
\usepackage[utf8]{inputenc}
\usepackage{lmodern}
\usepackage{amsmath,amssymb,amsthm,mathtools}
\usepackage{booktabs}
\usepackage{microtype}
\usepackage{hyperref}
\usepackage{enumitem}
\usepackage{tikz-cd}

\newtheorem{theorem}{Theorem}[section]
\newtheorem{lemma}[theorem]{Lemma}
\newtheorem{proposition}[theorem]{Proposition}
\newtheorem{corollary}[theorem]{Corollary}
\theoremstyle{definition}
\newtheorem{definition}[theorem]{Definition}
\theoremstyle{remark}
\newtheorem{remark}[theorem]{Remark}

\newcommand{\F}{\mathbb F}
\newcommand{\Q}{\mathbb Q}
\newcommand{\Z}{\mathbb Z}
\newcommand{\RA}{R_{\mathrm A}}
\newcommand{\Rstd}{R_{\mathrm{std}}}
\newcommand{\TQ}{T_{\mathbb Q}}
\newcommand{\wt}{\operatorname{wt}}
\newcommand{\Aut}{\operatorname{Aut}}
\newcommand{\Span}{\operatorname{span}}

\title{
Formal Root Reconstruction of $E_8$ from the Extended Binary Hamming Code\\
and Lean~4 Kernel Verification
}

\author{
Thomas Hoffbauer\\
\small Independent Research, Bamberg, Germany
}

\date{August 2026}

\begin{document}

\maketitle

% ============================================================================
% ABSTRACT
% ============================================================================

\begin{abstract}
We give an explicit and formally auditable construction relating the extended
binary Hamming code $C_8=[8,4,4]$ to the standard $240$-root realization of
the exceptional root system $E_8$.

Starting from the Fano-plane realization of the binary Simplex/Hamming pair
$[7,3,4]\subset[7,4,3]$, we distinguish two algebraically separate
construction branches.  The first yields the $[[7,1,3]]$ Steane stabilizer
code through a binary symplectic CSS construction.  The second extends the
Hamming code by a parity coordinate to the doubly-even self-dual code
$C_8=[8,4,4]$ and applies Construction~A.

The unscaled Construction~A minimal-vector set $\RA$ contains exactly
\[
16+224=240
\]
vectors of squared norm $4$.  We define the rational linear transformation
\[
\TQ=\frac12 H_2^{\oplus4},
\qquad
H_2=
\begin{pmatrix}
1&1\\
1&-1
\end{pmatrix},
\]
prove
\[
\TQ^{\,T}\TQ=\frac12I_8,
\]
construct an explicit left inverse, and establish injectivity.  The image of
$\RA$ is then shown to lie in the standard $E_8$ root set
\[
\Rstd=R_{\mathrm{int}}\cup R_{\mathrm{half}},
\qquad
|R_{\mathrm{int}}|=112,\quad |R_{\mathrm{half}}|=128.
\]
Injectivity together with the equality of finite cardinalities yields the
root-set identity
\[
\TQ(\RA)=\Rstd.
\]

The finite algebraic certificates are implemented in Lean~4/mathlib.
Verification status is attached directly to mathematical theorems rather
than Boolean status flags.  Classical runtime programs are used only as
independent companion witnesses.  No claim concerning physical quantum
hardware or quantum computational speedup is made.
\end{abstract}

% ============================================================================
\section{Introduction}
% ============================================================================

Relations among finite geometry, coding theory, quantum stabilizer codes, and
exceptional lattices are classical.  However, numerical coincidences among
quantities such as $7$, $8$, or $240$ do not by themselves establish a
structural connection.

The purpose of this work is therefore not to claim a new discovery of the
classical Fano--Hamming--Steane or Hamming--$E_8$ correspondences.  Rather,
we construct an explicit end-to-end chain whose individual algebraic maps,
finite enumerations, and root-set identification can be audited in a single
formal environment.

The main chain studied here is
\[
\text{Fano incidence}
\longrightarrow
[7,3,4]\subset[7,4,3]
\longrightarrow
[8,4,4]
\longrightarrow
\RA
\overset{\TQ}{\longrightarrow}
\Phi(E_8).
\]

A parallel branch produces the $[[7,1,3]]$ Steane code from the same
seven-coordinate Fano/Hamming data.

These two branches are deliberately kept separate:
\[
S_{\mathrm{Steane}}\subset\F_2^{14},
\qquad
\Rstd\subset\Q^8.
\]
There is no direct map from the stabilizer space to the $E_8$ root set in
the present construction.

\subsection*{Claim wall}

The results proved in this paper are statements in finite geometry, coding
theory, symplectic linear algebra, and lattice/root-system geometry.

In particular, we make no claims about:
\begin{itemize}[nosep]
    \item physical quantum hardware,
    \item decoherence or error rates of physical qubits,
    \item implementation of Shor order finding,
    \item quantum computational advantage,
    \item a physical realization of $E_8$.
\end{itemize}

% ============================================================================
\section{Fano geometry and the Steane branch}
% ============================================================================

\begin{definition}[Fano points]
We identify the seven points of the Fano plane with
\[
\F_2^3\setminus\{0\}.
\]
For distinct nonzero $x,y\in\F_2^3$, the corresponding line is
\[
\{x,y,x+y\}.
\]
\end{definition}

Every Fano line has three points.  Its complement therefore has four points.
The incidence vectors of the seven complements form the nonzero weight-$4$
words of the binary Simplex code $[7,3,4]$.

Let
\[
h_1,h_2,h_3\in\F_2^7
\]
be three independent complement vectors.  Then
\[
h_i\cdot h_j=0
\qquad
\text{for all }i,j.
\]

\subsection{Binary symplectic lift}

For $u=(x\mid z)$ and $v=(x'\mid z')$ in $\F_2^{14}$, define
\[
\omega(u,v)
=
x\cdot z'+z\cdot x'.
\]

This is the standard binary symplectic form underlying stabilizer
representations.

For each $h_i$ define
\[
X_i=(h_i\mid0),
\qquad
Z_i=(0\mid h_i).
\]

Since
\[
h_i\cdot h_j=0,
\]
all six generators are mutually symplectically orthogonal.

\begin{proposition}[Steane stabilizer isotropy]
Let
\[
S_{\mathrm{Steane}}
=
\Span_{\F_2}
\{X_1,X_2,X_3,Z_1,Z_2,Z_3\}.
\]
Then
\[
S_{\mathrm{Steane}}
\subseteq
S_{\mathrm{Steane}}^\perp
\]
and
\[
\dim S_{\mathrm{Steane}}=6.
\]
Consequently the associated stabilizer code has
\[
k=7-6=1
\]
logical qubit.
\end{proposition}

The minimum Pauli weight of
\[
S_{\mathrm{Steane}}^\perp
\setminus
S_{\mathrm{Steane}}
\]
is $3$, yielding the familiar parameters
\[
[[7,1,3]].
\]

The binary symplectic representation is precisely the type of structure used
in the stabilizer tableau formalism: Pauli commutation and anticommutation
are encoded by a symplectic product over $\F_2$.

% ============================================================================
\section{The extended Hamming code}
% ============================================================================

Let
\[
C_7=[7,4,3]
\]
denote the binary Hamming code dual to the Simplex code.

\begin{definition}[Parity extension]
For
\[
v=(v_0,\ldots,v_6)\in C_7,
\]
define
\[
\operatorname{parityLift}(v)
=
\left(
v_0,\ldots,v_6,
\sum_{i=0}^{6}v_i
\right).
\]
The sum is taken in $\F_2$.
\end{definition}

Its image is the extended Hamming code
\[
C_8=[8,4,4].
\]

\begin{theorem}[Weight enumerator]
The code $C_8$ has weight enumerator
\[
W_{C_8}(X,Y)
=
X^8+14X^4Y^4+Y^8.
\]
Hence the $16$ codewords consist of
\[
1\text{ word of weight }0,\qquad
14\text{ words of weight }4,\qquad
1\text{ word of weight }8.
\]
\end{theorem}

\begin{corollary}
The code $C_8$ has minimum distance $4$, is doubly even,
\[
\wt(c)\equiv0\pmod4
\qquad(c\in C_8),
\]
and is self-dual,
\[
C_8=C_8^\perp.
\]
\end{corollary}

% ============================================================================
\section{Construction A and its 240 minimal vectors}
% ============================================================================

\begin{definition}[Unscaled Construction A preimage]
Define
\[
L'_C
=
\left\{
x\in\Z^8:
x\bmod2\in C_8
\right\}.
\]
The usual Construction~A lattice is
\[
L(C_8)
=
\frac1{\sqrt2}L'_C.
\]
\end{definition}

We focus on the norm-$4$ vectors of the unscaled model:
\[
\RA
=
\{x\in L'_C:\|x\|^2=4\}.
\]

They split naturally into two sectors.

\subsection{Zero-codeword sector}

The zero codeword produces
\[
\pm2e_i,
\qquad i=1,\ldots,8.
\]
Thus
\[
|R_0|=8\cdot2=16.
\]

\subsection{Weight-four sector}

Each of the $14$ weight-$4$ codewords has a support of cardinality $4$.
Assigning independent signs on that support gives
\[
2^4=16
\]
minimal vectors per codeword.

Hence
\[
|R_4|
=
14\cdot16
=
224.
\]

The sectors are disjoint, so
\[
\boxed{
|\RA|
=
16+224
=
240.
}
\]

After the usual Construction~A scaling by $1/\sqrt2$, their norm becomes
\[
\frac12\cdot4=2,
\]
the standard squared length of $E_8$ roots.

% ============================================================================
\section{Rational root reconstruction}
% ============================================================================

Define
\[
H_2=
\begin{pmatrix}
1&1\\
1&-1
\end{pmatrix}
\]
and
\[
M=
H_2\oplus H_2\oplus H_2\oplus H_2.
\]

Explicitly,
\[
M=
\begin{pmatrix}
1&1&0&0&0&0&0&0\\
1&-1&0&0&0&0&0&0\\
0&0&1&1&0&0&0&0\\
0&0&1&-1&0&0&0&0\\
0&0&0&0&1&1&0&0\\
0&0&0&0&1&-1&0&0\\
0&0&0&0&0&0&1&1\\
0&0&0&0&0&0&1&-1
\end{pmatrix}.
\]

Set
\[
\TQ=\frac12M.
\]

\subsection{Gram identity}

Since
\[
H_2^TH_2=2I_2,
\]
one obtains
\[
M^TM=2I_8.
\]

Therefore
\[
\boxed{
\TQ^T\TQ
=
\frac12I_8.
}
\]

Consequently
\[
\|\TQ v\|^2
=
\frac12\|v\|^2.
\]

For every $v\in\RA$,
\[
\|v\|^2=4,
\]
hence
\[
\|\TQ v\|^2=2.
\]

\subsection{Explicit left inverse}

Define
\[
L=2\TQ^T.
\]

Then
\[
L\TQ
=
2\TQ^T\TQ
=
I_8.
\]

\begin{lemma}[Injectivity]
The linear map
\[
\TQ:\Q^8\to\Q^8
\]
is injective.
\end{lemma}

\begin{proof}
If $\TQ u=\TQ v$, then applying $L$ gives
\[
u
=
L\TQ u
=
L\TQ v
=
v.
\]
\end{proof}

% ============================================================================
\section{The standard $E_8$ root system}
% ============================================================================

We use the standard rational realization
\[
\Rstd
=
R_{\mathrm{int}}
\cup
R_{\mathrm{half}}.
\]

The integer sector consists of
\[
(\pm1,\pm1,0^6)
\]
with arbitrary positions of the two nonzero entries.  Thus
\[
|R_{\mathrm{int}}|
=
\binom82\cdot4
=
112.
\]

The half-integer sector consists of
\[
\frac12(\pm1,\ldots,\pm1)
\]
with an even number of minus signs.  Thus
\[
|R_{\mathrm{half}}|
=
2^7
=
128.
\]

Hence
\[
\boxed{
|\Rstd|=112+128=240.
}
\]

% ============================================================================
\section{Gate G.13: root-set image theorem}
% ============================================================================

The remaining structural task is to identify the image of the two
Construction~A sectors.

\begin{lemma}[Zero-sector containment]
\[
\TQ(R_0)\subseteq R_{\mathrm{int}}.
\]
\end{lemma}

Indeed, every vector $\pm2e_i$ is mapped to a vector containing exactly two
nonzero coordinates, each equal to $\pm1$.

\begin{lemma}[Weight-four-sector containment]
\[
\TQ(R_4)
\subseteq
\Rstd.
\]
\end{lemma}

This finite statement is checked sectorially in the formal implementation.
The action of each $2\times2$ Hadamard block is
\[
(a,b)
\longmapsto
\left(
\frac{a+b}{2},
\frac{a-b}{2}
\right),
\]
so weight-$4$ signed supports map either to integer roots or to the
half-integer parity sector.

Combining the two lemmas gives
\[
\boxed{
\TQ(\RA)\subseteq\Rstd.
}
\]

We can now state the main result.

\begin{theorem}[Exact root-set reconstruction]
\label{thm:root-reconstruction}
The rational transformation
\[
\TQ=\frac12H_2^{\oplus4}
\]
induces a bijection
\[
\boxed{
\TQ(\RA)=\Rstd.
}
\]
\end{theorem}

\begin{proof}
Injectivity gives
\[
|\TQ(\RA)|=|\RA|=240.
\]
By the sector containment theorem,
\[
\TQ(\RA)\subseteq\Rstd.
\]
But
\[
|\Rstd|=240.
\]
A finite subset of a finite set having the same cardinality is equal to the
whole set.  Therefore
\[
\TQ(\RA)=\Rstd.
\]
\end{proof}

% ============================================================================
\section{Lean 4 formalization}
% ============================================================================

The formal development separates mathematical certificates into small gates.
The purpose of this decomposition is epistemic as well as technical:
a global status is never represented merely by a Boolean flag.

\begin{center}
\begin{tabular}{lll}
\toprule
Gate & Statement & Status\\
\midrule
G.13a &
$M^TM=2I_8$, $\TQ^T\TQ=\frac12I_8$
&
kernel theorem
\\
G.13b &
explicit left inverse and injectivity
&
kernel theorem
\\
G.13c &
$|\RA|=240$
&
kernel theorem
\\
G.13d &
$\TQ(\RA)\subseteq\Rstd$
&
kernel theorem
\\
G.13e &
$|\Rstd|=240$
&
kernel theorem
\\
G.13f &
$\TQ(\RA)=\Rstd$
&
derived kernel theorem
\\
\bottomrule
\end{tabular}
\end{center}

The final Lean argument follows the same finite-cardinality proof as
Theorem~\ref{thm:root-reconstruction}.  Schematically,

\begin{verbatim}
theorem g13_root_set_bijection :
    R_A.image TQ.mulVec = R_std := by
  apply Finset.eq_of_subset_of_card_le r_A_subset_r_std
  rw [Finset.card_image_of_injective _ tq_injective]
  rw [r_A_card_240, r_std_card_240]
\end{verbatim}

The repository governance principle is that the final status theorem is
logically coupled to this equality itself, rather than to an unrelated
Boolean marker.

% ============================================================================
\section{Independent runtime witnesses}
% ============================================================================

Finite C++ companion programs reproduce selected cardinalities, sector
decompositions, and arithmetic identities.

These programs are \emph{not} treated as formal proofs.  Their role is to
provide independent runtime witnesses and regression tests for the same
finite structures represented in Lean.

The epistemic hierarchy is therefore

\[
\boxed{
\text{runtime witness}
\neq
\text{kernel proof}.
}
\]

Agreement between the two is useful for software validation, but only the
formal theorem establishes the corresponding kernel-level mathematical
certificate.

% ============================================================================
\section{Relation to the Steane/Fano branch}
% ============================================================================

The $E_8$ reconstruction and the Steane stabilizer construction share the
same binary Fano/Hamming origin, but belong to different ambient categories.

The structural diagram is

\[
\begin{array}{ccccc}
&&\text{Fano incidence}&&\\[1mm]
&\swarrow&&\searrow&\\
[7,3,4]&&&&[7,4,3]\\
\downarrow&&&&\downarrow\\
[[7,1,3]]&&&&[8,4,4]\\
&&&&\downarrow\\
&&&&\RA\\
&&&&\downarrow\TQ\\
&&&&\Phi(E_8).
\end{array}
\]

This diagram should not be interpreted as a direct map from the quantum
stabilizer space into the $E_8$ root space.  Rather, the two constructions
share a common finite-geometric origin.

% ============================================================================
\section{Scope and limitations}
% ============================================================================

Several limitations are essential.

First, the individual correspondences among the Fano plane, Hamming codes,
the Steane code, Construction~A, and $E_8$ are classical.  The contribution
here is their explicit integration into one audited formal architecture.

Second, the theorem
\[
\TQ(\RA)=\Rstd
\]
is a root-set identification.  It should not be silently strengthened to a
formal equality of complete lattice objects unless the corresponding lattice
generation theorem is also formalized.

Third, no statement about physical quantum hardware follows from the
stabilizer construction.

Finally, no computational speedup is claimed from the existence of a
classical stabilizer simulation or from the finite runtime companion
programs.

% ============================================================================
\section{Conclusion}
% ============================================================================

We have given an explicit root reconstruction beginning with the extended
binary Hamming code $[8,4,4]$.

Construction~A produces a norm-$4$ root set
\[
\RA=R_0\sqcup R_4
\]
with
\[
|R_0|=16,
\qquad
|R_4|=224,
\qquad
|\RA|=240.
\]

The rational map
\[
\TQ=\frac12H_2^{\oplus4}
\]
satisfies
\[
\TQ^T\TQ=\frac12I_8
\]
and possesses the explicit left inverse
\[
2\TQ^T.
\]

Its image is contained in the standard $E_8$ root set and has the same
cardinality.  Hence
\[
\boxed{
\TQ(\RA)=\Phi(E_8).
}
\]

The main value of the formal development is not the rediscovery of these
classical correspondences, but the construction of an end-to-end,
machine-auditable chain in which algebraic claims, finite computational
witnesses, and formal kernel certificates are kept explicitly distinct.

% ============================================================================
\begin{thebibliography}{99}

\bibitem{AaronsonGottesman}
S.~Aaronson and D.~Gottesman,
\newblock Improved Simulation of Stabilizer Circuits,
\newblock \emph{Physical Review A} \textbf{70}, 052328 (2004),
\newblock arXiv:quant-ph/0406196.

\bibitem{ConwaySloane}
J.~H. Conway and N.~J.~A. Sloane,
\newblock \emph{Sphere Packings, Lattices and Groups},
\newblock 3rd ed., Springer, New York, 1999.

\bibitem{Steane}
A.~M. Steane,
\newblock Error Correcting Codes in Quantum Theory,
\newblock \emph{Physical Review Letters} \textbf{77},
793--797 (1996).

\bibitem{GottesmanThesis}
D.~Gottesman,
\newblock \emph{Stabilizer Codes and Quantum Error Correction},
\newblock Ph.D. thesis, California Institute of Technology, 1997.

\bibitem{MacWilliamsSloane}
F.~J. MacWilliams and N.~J.~A. Sloane,
\newblock \emph{The Theory of Error-Correcting Codes},
\newblock North-Holland, Amsterdam, 1977.

\end{thebibliography}

\end{document}